\documentclass[11pt]{amsart}
\usepackage[margin=1in]{geometry}
\usepackage{amssymb,amsmath,amsthm,mathtools}
\usepackage{booktabs,array}
\usepackage[hidelinks]{hyperref}
\hypersetup{pdftitle={Spanning Trees per Edge: A Sharp Bound and a Question of Kenyon, with a Lean-Checked Key Inequality and Computer Certificates},pdfauthor={Filip Majchrzak}}
\IfFileExists{microtype.sty}{\usepackage[expansion=false]{microtype}}{}
\emergencystretch=3em

\theoremstyle{plain}
\newtheorem{theorem}{Theorem}
\newtheorem{proposition}[theorem]{Proposition}
\newtheorem{lemma}[theorem]{Lemma}
\newtheorem{corollary}[theorem]{Corollary}
\theoremstyle{remark}
\newtheorem{remark}[theorem]{Remark}

\newcommand{\R}{\mathbb{R}}
\newcommand{\cC}{\mathcal{C}}
\newcommand{\cG}{\mathcal{G}}
\newcommand{\code}[1]{\texttt{#1}}
\newcommand{\rb}{\bar r}
\newcommand{\tht}{\theta}
\DeclareMathOperator{\tr}{tr}

\title[Spanning trees per edge]{Spanning Trees per Edge:\\ A Sharp Bound and a Question of Kenyon\\[4pt] {\normalfont\large with a Lean-checked key inequality and computer certificates}}
\author{Filip Majchrzak}
\address{Independent researcher, Poland (SAPIO MORPHEUS project)}
\thanks{ORCID 0009-0007-3475-8900. Version 1.0, 24 November 2026. Archived at Zenodo, \href{https://doi.org/10.5281/zenodo.24504455}{doi:10.5281/zenodo.24504455}. Text licensed CC BY 4.0; code Apache-2.0.}
\date{24 November 2026}

\begin{document}

\begin{abstract}
Let $\tau(G)$ be the number of spanning trees of a finite multigraph $G$ with $m$ edges. We prove that
\[
  \tau(G)\;<\;\Bigl(\tfrac{27}{8}\Bigr)^{m/2}\;=\;(1.8371\ldots)^m\qquad(m\ge1),
\]
and that the base $\sqrt{27/8}$ cannot be lowered: it is approached by $4$-regular graphs of large girth, by
theorems of McKay and Lyons. In 1996 Kenyon asked for the supremum of $\tau(G)^{1/m}$ over all graphs and over planar graphs;
the first supremum is therefore $\sqrt{27/8}$. For planar graphs, equivalently for determinants of link diagrams,
the bound improves the constant $1.8393$ of Stoimenow slightly, to $1.8371$; the conjectured planar value
$1.7916$ remains open, and no argument that ignores planarity can go below $\sqrt{27/8}$. The proof compares
$G$ with its universal cover, uses a known closed formula for the tree entropy of the cover, and bounds it edge
by edge through one concavity step and an explicit one-variable inequality.
The a priori bounds, the two identities and the key inequality (Sections \ref{sec:cover} and \ref{sec:edge}) are formalised in Lean~4 with Mathlib as statements about real numbers; the probabilistic input, the closed formula, the comparison with the cover and the reductions are ordinary proofs and are not formalised, so the theorem as a whole is not machine-checked. The numerical inequalities are also certified by interval arithmetic, and the theorem was confirmed exhaustively for all multigraphs with at most $14$ edges.
\end{abstract}

\maketitle

\section{Introduction}

All graphs are finite multigraphs: parallel edges and loops are allowed. We write $m$ for the number of edges and
$\tau(G)$ for the number of spanning trees ($\tau=0$ if $G$ is disconnected). Throughout,
\[
  \tht:=\sqrt{27/8}=1.8371173\ldots,\qquad L:=\log\tht=\tfrac12\log\tfrac{27}{8}=0.6081976\ldots
\]

\begin{theorem}\label{thm:main}
Every multigraph $G$ with $m\ge1$ edges satisfies $\tau(G)<\tht^{\,m}=(27/8)^{m/2}$.
\end{theorem}

\begin{proposition}[McKay; Lyons]\label{prop:sharp}
If $G_j$ are connected $4$-regular graphs whose girth tends to infinity, then $\tau(G_j)^{1/m(G_j)}\to\tht$.
\end{proposition}

\begin{corollary}\label{cor:sup}
$\sup_G\tau(G)^{1/m}=\sqrt{27/8}$, the supremum being over all multigraphs with at least one edge; it is not attained.
\end{corollary}

Proposition~\ref{prop:sharp} is classical: McKay \cite{McKay} determined the constant $(d-1)^{d-1}/(d^2-2d)^{d/2-1}$
for $d$-regular graphs, which is $27/8$ per vertex for $d=4$, and Lyons \cite{Lyons} (Theorem~3.3 and
Example~3.16 in the arXiv version) proved the convergence for every sequence of connected graphs whose random weak
limit is the $4$-regular tree; girth tending to infinity implies this. A $4$-regular graph with $N$ vertices has
$m=2N$ edges. Corollary~\ref{cor:sup} follows from Theorem~\ref{thm:main} and Proposition~\ref{prop:sharp}.

\subsection*{A question of Kenyon}
In his paper on tiling a rectangle with the fewest squares, Kenyon \cite{Kenyon} (\S9 of the arXiv version) needed
an exponential bound
for the number of spanning trees of a planar graph and wrote:
\begin{quote}
``\textbf{Problem 2.} Among all graphs with $m$ edges, which graph $G$ has the largest number of spanning trees
$\kappa(G)$? What is the sup of $\kappa(G)^{1/m}$ over all graphs? Over planar graphs?''
\end{quote}
He records that the true value lies in the range $(1.79,2)$: the trivial bound is $2^m$, he mentions an unpublished
bound $\lambda^m$ with some $\lambda<2$ due to N.~Young, and the square grid gives $e^{2C/\pi}=1.7916\ldots$ ($C$ is
Catalan's constant). Corollary~\ref{cor:sup} answers the second of his three questions: \emph{the supremum over
all graphs is $\sqrt{27/8}$.} The first question (the extremal graph for each $m$) and the third (planar graphs) are not
answered here.

\subsection*{Planar graphs and link determinants}
For planar graphs the conjectured value is $e^{2C/\pi}$; it is attributed to Kenyon in
\cite[Conjecture~2.3]{CKP}, where it is shown to be equivalent to a conjecture on determinants of links. The
authors of \cite{CKP} also record, from Kenyon, that the conjectured planar bound fails for some non-planar
graphs; Proposition~\ref{prop:sharp} is a quantitative form of this. Proven upper bounds for planar graphs are
$1.8637^m$ (Fomin \cite{Fomin}) and $\delta^m$ with $\delta\approx1.8393$ the real root of $x^3=x^2+x+1$
(Stoimenow \cite{Stoimenow}, stated for link diagrams; we quote it from \cite{CKP,GJKLZ}, see
Section~\ref{sec:lit}). We have not found a better one in the literature.

\begin{corollary}\label{cor:planar}
\textup{(a)} Every planar multigraph with $m\ge1$ edges has fewer than $\tht^{\,m}$ spanning trees, so the planar
supremum in Kenyon's problem lies in $[\,e^{2C/\pi},\sqrt{27/8}\,]=[1.7916\ldots,1.8371\ldots]$.

\textup{(b)} For every link diagram $K$ with $k\ge1$ crossings, $\det(K)<(27/8)^{k/2}$; in the normalisation of
\cite{CKP}, $2\pi\log\det(K)/k<\pi\log(27/8)=3.8214\ldots$
\end{corollary}
\noindent The corresponding numbers are $3.8288$ for Stoimenow's bound and $v_{\mathrm{oct}}=3.6639$ for the
conjecture, where $v_{\mathrm{oct}}$ is the volume of the regular ideal hyperbolic octahedron. Part~(b) follows
from~(a): the Tait graph of $K$ is planar with $k$ edges, and $\det(K)$ is a signed count of its spanning trees
\cite[Lemma~2.1]{CKP}, hence at most $\tau$. The improvement over \cite{Stoimenow} is
small. Its interest is that $\tht$ is exact for general graphs, so it marks the limit of every method that does
not use planarity.

\begin{corollary}\label{cor:current}
Let $\Lambda_m$ be the largest value of $1/|I|$ over the nonzero branch currents $I$ of planar networks of $m$ unit
resistors, with a unit current fed through two terminals on a common face, and let $c=\lim\Lambda_m^{1/m}$ be the
growth constant studied in \cite{M2}. Then $1.7194145\le c\le\tht=1.8371173\ldots$
\end{corollary}
\noindent Indeed, under a unit current feed every nonzero branch current of a network $G$ has the form $N/\tau(G)$
with an integer $N\ne0$ (Kirchhoff), so its reciprocal is at most $\tau(G)<\tht^{\,m}$; the lower bound is
from~\cite{M2}. This replaces the upper bound $1.8393$ used there.

\subsection*{Method, and what is new}
The proof has three parts.
\begin{enumerate}
\item \emph{Comparison with the universal cover.} Return probabilities of the random walk on $G$ dominate those on
  its universal cover $T$, which gives $\log\tau(G)\le\Phi(G)+O(1)$, where $\Phi(G)$ is the tree entropy of $T$ per
  fundamental domain. This is the method of McKay \cite{McKay} and Lyons \cite{Lyons}.
\item \emph{A closed formula.} $\Phi(G)=\sum_e\log(1+r_e)-\sum_v\log x_v$, where the positive numbers $x_v$ and
  $r_e$, defined in Section~\ref{sec:cover} from hitting probabilities on $T$, satisfy $r_e(1+r_e)=x_ux_v$ and
  $\sum_{e\ni v}r_e=d_vx_v-1$. This is known: it is the Floquet formula of
  Banks, Breuer, Garza-Vargas, Seelig and Simon \cite{BBGSS} at energy $0$, and the system is the covariance system
  of Gaussian belief propagation, see Watanabe \cite{Watanabe}. We include a short proof.
\item \emph{An edge inequality.} We show $\Phi(G)\le mL$ with equality exactly for $4$-regular graphs. A
  multiplier $\tfrac18$ on the edge equations makes the vertex terms concave, and Jensen's inequality reduces
  everything to the one-variable functions
  $f_d(r)=\tfrac9{16}\log(1+r)+\tfrac1{16}\log r-(\tfrac1d+\tfrac18)\log(r+\tfrac1d)$, whose derivative has the sign
  of $(8r-1)((d-2)r-1)$ and whose value at $d=4$, $r=\tfrac18$ is exactly $\tfrac14\log\tfrac{27}{8}$.
\end{enumerate}
Parts 1 and 2 are not new. As far as we could determine, part~3 and Theorem~\ref{thm:main} are. The literature
search and its limits are described in Section~\ref{sec:lit}; we make no claim beyond it.

\subsection*{Status of verification}
The proof is elementary but has several steps, and it was produced by AI agents under the author's direction (see
the end of the paper). We therefore state exactly what has been checked and how.
\begin{itemize}
\item \emph{Formal verification of the key inequality (Lean~4 with Mathlib).} Sections \ref{sec:cover} and
  \ref{sec:edge} are formalised as statements about real numbers. Lean checks the following: let $h$ be a positive
  function on the darts of a finite multigraph with minimum degree at least $2$ and no edge between two vertices of
  degree $2$; if $h$ satisfies the first-step equation \eqref{eq:h} and is the pointwise limit of functions
  $h^{(t)}$ with $h^{(0)}=0$ that satisfy the inequality stated after \eqref{eq:h}, then
  $\sum_v\log\bigl(d_v-\sum_{u\sim v}h_{u\to v}\bigr)-\sum_e\log(1-q_e)\le m\cdot\frac12\log\frac{27}{8}$; the left
  side equals $\sum_e\log(1+r_e)-\sum_v\log x_v$. This covers the bounds of Lemma~\ref{lem:apriori}, the
  identities \eqref{eq:I}, the bounds on $\rb_v$ in Lemma~\ref{lem:ranges}, Lemma~\ref{lem:onevar} and the
  inequality $E_e\le L$ of Lemma~\ref{lem:edge}. In place of the decimal values of Lemma~\ref{lem:num} and of
  Case~C, Lean uses exact inequalities between natural numbers for degrees $2$ to $119$, evaluated by its kernel,
  and a direct estimate for larger degrees. The proof contains no \code{sorry} and uses only the three standard
  axioms of Lean and Mathlib, and an explicit example shows that equality can occur.
  \emph{Not formalised:} that hitting probabilities on the universal cover satisfy these first-step relations, the
  closed formula (Lemma~\ref{lem:formula}), Lemmas \ref{lem:mtt} and \ref{lem:compare}, the reductions of
  Section~\ref{sec:red}, the amplification of Section~\ref{sec:proof} and the equality statement of
  Remark~\ref{rem:eq}. Theorem~\ref{thm:main} itself is therefore not machine-checked.
\item \emph{Interval arithmetic.} Independently of Lean, the numerical facts of Lemma~\ref{lem:num} and the
  constants of Case~C of Lemma~\ref{lem:edge} are certified with $60$-digit interval arithmetic at explicit rational
  points (Appendix~\ref{app:num}).
\item \emph{Recomputation by separate agents.} Separate AI agents of the same project re-derived each step of the
  proof and of this text and recomputed the numerical claims with their own code.
\item \emph{Exhaustive check.} Theorem~\ref{thm:main} was confirmed directly, in exact integer arithmetic, for all
  $8{,}350{,}006$ connected loopless multigraphs with at most $14$ edges that are not trees (for a tree $\tau=1$).
  This is a sanity check far from the sharp regime.
\end{itemize}
What has not been done: no external expert has reviewed this paper. We welcome corrections.

\section{Reductions}\label{sec:red}

We prove $\tau(G)\le\tht^{\,m}$ for all $m\ge0$ by induction on $m$. This implies Theorem~\ref{thm:main}, because
for $m\ge1$ the number $(27/8)^{m/2}$ is not an integer (for even $m$ it is a non-integral rational, for odd $m$
it is irrational). The case $m=0$ is trivial, and disconnected graphs have $\tau=0$. Let $G$ be connected with
$m\ge1$ edges, assume the claim for fewer edges, and apply the first of the following cases that fits.

\begin{itemize}
\item[(R1)] $G$ has a loop. Deleting it keeps $\tau$ and lowers $m$.
\item[(R2)] $G$ has a vertex of degree $1$. Deleting it keeps $\tau$ and lowers $m$.
\item[(R3)] $G$ is a cycle with $k\ge2$ edges. Then $\tau=k\le\tht^{\,k}$.
\item[(R4)] $G$ has two adjacent vertices $x,y$ of degree $2$. Let $W$ be the maximal chain through the edge $xy$
  all of whose interior vertices have degree $2$. Since (R1)--(R3) do not apply, $W$ has $k\ge3$ edges and its end
  vertices $u,v$ have degree at least $3$.
  If $u=v$, then $W$ is a cycle attached at $u$; deleting its interior vertices leaves a graph $G''$ with $m-k$
  edges and $\tau(G)=k\,\tau(G'')\le k\,\tht^{\,m-k}\le\tht^{\,m}$.
  If $u\ne v$, let $G'$ be $G$ with $W$ replaced by a single edge $e=uv$, so $m(G')=m-k+1$. A spanning tree of $G$
  contains all of $W$, or all edges of $W$ but one; this gives
  $\tau(G)=\tau_e(G')+k\,\tau_{\bar e}(G')\le k\,\tau(G')$, where $\tau_e$ and $\tau_{\bar e}$ count the spanning trees
  of $G'$ with and without $e$. Hence $\tau(G)\le k\,\tht^{\,m-k+1}\le\tht^{\,m}$, because $k\le\tht^{\,k-1}$ for
  $k\ge3$ (indeed $3\le27/8$ and $(k+1)/k\le4/3<\tht$).
\item[(R5)] Otherwise $G$ belongs to the class $\cC$ of \emph{connected loopless multigraphs with at least one
  edge, minimum degree at least $2$, and no two adjacent vertices of degree $2$.}
\end{itemize}
Note that $k=2$ would fail in (R4), since $2>\tht$: a vertex of degree $2$ without a neighbour of degree $2$
cannot be removed, and such vertices stay in $\cC$. Every $G\in\cC$ has a vertex of degree at least $3$, and gluing two members of $\cC$ at a vertex gives a
member of $\cC$. The rest of the paper proves $\tau(G)\le\tht^{\,m}$ for $G\in\cC$.

\section{The universal cover}\label{sec:cover}

Let $G\in\cC$ with vertex degrees $d_v$, let $T$ be its universal cover and consider simple random walk on the
tree $T$. A \emph{dart} is an edge with a direction, written $u\to v$. For adjacent vertices $u,v$ of $T$ let
$h_{u\to v}$ be the probability that the walk started at $u$ ever visits $v$. All quantities below are invariant under deck transformations and are regarded as functions of
the darts, vertices and edges of $G$. Conditioning on the first step, and using that in a tree a walk from a
neighbour $w\ne v$ of $u$ must return to $u$ before reaching $v$,
\begin{equation}\label{eq:h}
 h_{u\to v}=\frac1{d_u}\Bigl(1+\sum_{w\sim u,\,w\ne v}h_{w\to u}\,h_{u\to v}\Bigr),\qquad\text{i.e.}\qquad
 h_{u\to v}=\frac1{d_u-\sum_{w\ne v}h_{w\to u}} .
\end{equation}
The sum is over the $d_u-1$ other darts into $u$. For the probabilities $h^{(t)}$ of a visit within $t$ steps the
same argument gives the inequality
$h^{(t+1)}_{u\to v}\le\frac1{d_u}\bigl(1+\sum_{w\ne v}h^{(t)}_{w\to u}h^{(t+1)}_{u\to v}\bigr)$.

\begin{lemma}\label{lem:apriori}
For $G\in\cC$: $h_{w\to u}\le\frac34$ if $d_w=2$, and $h_{u\to v}\le\frac4{d_u+3}\le\frac23$ if $d_u\ge3$. In
particular the walk on $T$ is transient.
\end{lemma}
\begin{proof}
Induction on $t$ for $h^{(t)}$, with $h^{(0)}=0$, carrying both bounds for all darts. If $d_w=2$, the other
neighbour of $w$ has degree at least $3$, so the incoming value is at most $\frac23$ and
$h^{(t+1)}_{w\to u}\le1/(2-\frac23)=\frac34$. If $d_u=d\ge3$, all incoming values are at most $\frac34$, so
$h^{(t+1)}_{u\to v}\le1/(d-\frac34(d-1))=\frac4{d+3}$. Let $t\to\infty$.
\end{proof}

For a vertex $v$ and an edge $e=uv$ put
\[
 x_v:=\frac1{d_v-\sum_{u\sim v}h_{u\to v}},\qquad q_e:=h_{u\to v}h_{v\to u}\in(0,1),\qquad
 r_e:=\frac{q_e}{1-q_e}>0,\qquad \rb_v:=\frac1{d_v}\sum_{e\ni v}r_e .
\]
Thus $d_vx_v$ is the Green function $\sum_kp_k(v,v)$ of the walk on $T$, where $p_k(v,v)$ denotes the $k$-step
return probability.

\begin{lemma}\label{lem:ident}
For every edge $e=uv$ and every vertex $v$,
\begin{equation}\label{eq:I}
 r_e(1+r_e)=x_ux_v,\qquad\qquad \sum_{e\ni v}r_e=d_vx_v-1,\ \text{ i.e. }\ \rb_v=x_v-\tfrac1{d_v}.
\end{equation}
\end{lemma}
\begin{proof}
By \eqref{eq:h}, $1/h_{v\to u}=1/x_v+h_{u\to v}$, hence $h_{v\to u}=(1-q_e)x_v$ and likewise
$h_{u\to v}=(1-q_e)x_u$. Multiplying gives $q_e=(1-q_e)^2x_ux_v$, the first identity. For the second,
$d_vx_v-1=x_v\sum_{u\sim v}h_{u\to v}=\sum_{e\ni v}(1-q_e)x_ux_v=\sum_{e\ni v}r_e$.
\end{proof}

\begin{lemma}\label{lem:ranges}
Let $G\in\cC$ and let $v$ be a vertex of degree $d$.
\begin{itemize}
\item[(a)] If $d=2$, then $0<\rb_v\le1$ and $x_v\le\frac32$.
\item[(b)] If $d\ge3$, then $r_e\le\frac3d$ for every edge at $v$; hence $0<\rb_v\le\frac3d$ and $x_v\le\frac4d$.
\item[(c)] If $d\ge3$ and $v$ has a neighbour of degree at least $3$, then
 $\rb_v\le R^*_d:=\dfrac{9d-1}{d(3d+1)}$.
\end{itemize}
\end{lemma}
\begin{proof}
(a) Both neighbours have degree at least $3$, so $x_v\le1/(2-\frac43)=\frac32$ and $\rb_v=x_v-\frac12\le1$.
(b) By Lemma~\ref{lem:apriori}, $q_e\le\frac34\cdot\frac4{d+3}=\frac3{d+3}$, so $r_e\le\frac3d$.
(c) One of the darts into $v$ starts at a vertex of degree at least $3$ and is at most $\frac23$; the others are at
most $\frac34$. Hence $\sum_{u\sim v}h_{u\to v}\le\frac34(d-1)+\frac23=\frac34d-\frac1{12}$, so
$x_v\le\frac{12}{3d+1}$ and $\rb_v=x_v-\frac1d\le\frac{9d-1}{d(3d+1)}$.
\end{proof}
For $d=3$ the bounds of Lemmas \ref{lem:apriori} and \ref{lem:ranges}(a),(b) are attained by the cover of the theta
graph with each edge subdivided once. The refinement (c) is needed only for $d=5$.

\section{A closed formula for the tree entropy of the cover}\label{sec:formula}

Define
\[
 \Phi(G):=\sum_{v\in V(G)}\Bigl(\log d_v-\sum_{k\ge1}\frac{p_k(v,v)}{k}\Bigr).
\]
By \cite{Lyons}, $\Phi(G)/|V(G)|$ is the limit of $\log\tau$ per vertex along finite connected graphs of bounded
degree whose random weak limit is $T$ rooted at the lift of a uniform vertex of $G$; we do not use this.

\begin{lemma}[known; see \cite{BBGSS,Watanabe}]\label{lem:formula}
$\displaystyle\Phi(G)=\sum_{e\in E(G)}\log(1+r_e)-\sum_{v\in V(G)}\log x_v$.
\end{lemma}
\begin{proof}
For $z\in[0,1]$ let $h_{u\to v}(z)$ be the generating function of the first hitting time of $v$ from $u$. As in
\eqref{eq:h}, $h_{u\to v}(z)=z/(d_u-z\sum_{w\ne v}h_{w\to u}(z))$. Put $P_v(z):=d_v-z\sum_{u\sim v}h_{u\to v}(z)$ and
$q_e(z):=h_{u\to v}(z)h_{v\to u}(z)$. Then $\cG_v(z):=\sum_kp_k(v,v)z^k=d_v/P_v(z)$, and the recursion gives, for
$z>0$,
\begin{equation}\label{eq:star}
 P_v(z)=z\,\frac{1-q_e(z)}{h_{v\to u}(z)}\qquad\text{for every edge }e=uv .
\end{equation}
Let $\Xi(z):=\sum_v\log P_v(z)-\sum_e\log(1-q_e(z))$. In $\Xi'(z)$ the coefficient of $h'_{u\to v}(z)$ is
$-z/P_v(z)+h_{v\to u}(z)/(1-q_e(z))$, which vanishes by \eqref{eq:star}. Therefore
\[
 \Xi'(z)=-\sum_v\frac{\sum_{u\sim v}h_{u\to v}(z)}{P_v(z)}=-\sum_v\frac{\cG_v(z)-1}{z}.
\]
By the paragraph ``Uniform decay'' in the proof of Lemma~\ref{lem:compare}, which uses only the first-step
inequality, all $h_{u\to v}(z)$ are analytic for $|z|<\frac{101}{100}$, and $P_v(z)>0$, $q_e(z)<1$ on $[0,1]$; so
$\Xi$ is analytic on $[0,1]$. Moreover $\Xi(0)=\sum_v\log d_v$. Integrating, $\Xi(1)=\sum_v\log d_v-\sum_v\int_0^1(\cG_v(z)-1)\,dz/z=\Phi(G)$. Finally
$P_v(1)=1/x_v$ and $1/(1-q_e)=1+r_e$.
\end{proof}

\section{Comparison of a graph with its cover}\label{sec:compare}

\begin{lemma}\label{lem:mtt}
Let $G$ be a connected loopless multigraph with $n\ge2$ vertices, adjacency matrix $A$ and diagonal degree
matrix $\Delta$, and let $1=\lambda_1>\lambda_2\ge\dots\ge\lambda_n\ge-1$ be the eigenvalues of the walk matrix
$M=\Delta^{-1}A$. Then
$\tau(G)=\dfrac{\prod_vd_v}{2m}\prod_{i\ge2}(1-\lambda_i)$.
\end{lemma}
\begin{proof}
The symmetric matrix $\mathcal L=I-\Delta^{-1/2}A\Delta^{-1/2}$ has eigenvalues $1-\lambda_i$ and a one-dimensional kernel
spanned by the unit vector $\varphi_v=\sqrt{d_v/2m}$, so its adjugate is $\prod_{i\ge2}(1-\lambda_i)\,\varphi\varphi^{\mathsf T}$.
The $(v,v)$ entry of the adjugate is the corresponding principal minor of $\mathcal L$, which equals the minor of
the Laplacian $\Delta-A$, namely $\tau(G)$ by the matrix-tree theorem, divided by $\prod_{u\ne v}d_u$.
\end{proof}

\begin{lemma}\label{lem:compare}
There is an absolute constant $C_1$ such that $\log\tau(G)\le\Phi(G)+C_1$ for every $G\in\cC$.
\end{lemma}
\begin{proof}
For $\lambda\in[-1,1)$ and $K\ge1$ we have $\log(1-\lambda)\le-\sum_{k\le K}\lambda^k/k+\frac1{K+1}$: for
$\lambda\ge0$ the omitted terms are non-positive, and for $\lambda<0$ the tail is an alternating series with
decreasing terms. Summing over $i\ge2$ in Lemma~\ref{lem:mtt}, with $\sum_{i\ge2}\lambda_i^k=\tr M^k-1$,
\[
 \log\tau(G)\le\sum_v\log d_v-\log 2m-\sum_{k\le K}\frac{\tr M^k-1}{k}+\frac{n-1}{K+1}.
\]
Now $\tr M^k=\sum_vp^G_k(v,v)\ge\sum_vp_k(v,v)$, because a walk on $T$ that returns to its starting point projects
to a walk on $G$ that returns. Hence
\begin{equation}\label{eq:cmp}
 \log\tau(G)\le\Phi(G)+\sum_v\sum_{k>K}\frac{p_k(v,v)}{k}+(1+\log K)+\frac{n-1}{K+1}-\log 2m .
\end{equation}
\emph{Uniform decay.} Let $z_0=\frac{101}{100}$, $\alpha=\frac{77}{100}$, $\beta=\frac{83}{100}$. Then
$\alpha\le\beta$, $\beta\ge z_0/(2-z_0\alpha)$, $\alpha\ge z_0/(3-2z_0\beta)$ and $z_0\beta<1$ (exact rational
arithmetic). The induction of Lemma~\ref{lem:apriori}, applied to the truncated expectations
$E_u[z_0^{\,\sigma};\sigma\le t]$ of the hitting time $\sigma$ of $v$ (finite sums, for which the first-step
inequality holds, with positive denominators because $z_0\beta<1$), gives $h_{w\to u}(z_0)\le\beta$ if $d_w=2$ and
$h_{u\to v}(z_0)\le\alpha$ if $d_u\ge3$; here all incoming values are at most $\beta$, and
$z_0/(d-(d-1)z_0\beta)$ is largest at $d=3$. Consequently the first-return generating function satisfies
$\frac{z_0}{d_v}\sum_uh_{u\to v}(z_0)\le z_0\beta<1$, so $\cG_v(z_0)\le1/(1-z_0\beta)<7$ and
$p_k(v,v)\le7\,(100/101)^k$ for all $k$ and all $v$, uniformly over $\cC$.

Take $K=1000\,n$ in \eqref{eq:cmp}. The double sum is less than $1$, $(n-1)/(K+1)<1$, and
$\log K-\log 2m\le\log 500$ because $n\le m$ in $\cC$. So $C_1=3+\log 500$ works.
\end{proof}

\begin{remark}
The inequality $\tau(G)\le e^{\Phi(G)}$ is false in general: for the theta graph with each edge subdivided once,
$\tau=12$ and $\log\tau-\Phi=0.1178$. The additive constant is removed in Section~\ref{sec:proof} by an
amplification argument.
\end{remark}

\section{The edge inequality}\label{sec:edge}

Let $g(r):=\frac9{16}\log(1+r)+\frac1{16}\log r$ for $r>0$; $g$ is concave and increasing. For an edge $e=uv$ with
$a=d_u$, $b=d_v$ define
\[
 E_e:=\log(1+r_e)-\tfrac1a\log x_u-\tfrac1b\log x_v-g'(\rb_u)(r_e-\rb_u)-g'(\rb_v)(r_e-\rb_v).
\]

\begin{lemma}\label{lem:sum}
\textup{(a)} $\sum_{e}E_e=\Phi(G)$. \quad\textup{(b)} $E_e\le f_a(\rb_u)+f_b(\rb_v)$, where
\[
 f_d(r):=g(r)-\bigl(\tfrac1d+\tfrac18\bigr)\log\bigl(r+\tfrac1d\bigr).
\]
\end{lemma}
\begin{proof}
(a) Summing over the edges at a fixed vertex $v$, the terms $g'(\rb_v)(r_e-\rb_v)$ cancel by the definition of
$\rb_v$, and $\sum_e(\frac1a\log x_u+\frac1b\log x_v)=\sum_v\log x_v$; use Lemma~\ref{lem:formula}.
(b) By \eqref{eq:I}, $\log(1+r_e)=2g(r_e)-\frac18(\log x_u+\log x_v)$. So $E_e$ is the sum over the two ends
$w\in\{u,v\}$ of $g(r_e)-g'(\rb_w)(r_e-\rb_w)-(\frac1{d_w}+\frac18)\log x_w$. By concavity,
$g(r_e)-g'(\rb_w)(r_e-\rb_w)\le g(\rb_w)$, and $x_w=\rb_w+\frac1{d_w}$.
\end{proof}

\begin{lemma}\label{lem:onevar}
For $d\ge2$ and $r>0$, $f_d'(r)$ has the sign of $Q_d(r):=(8r-1)\bigl((d-2)r-1\bigr)$. Hence $f_2$ increases on
$(0,\frac18]$ and decreases afterwards; for $d\ge3$, $f_d$ increases up to $\min(\frac18,\frac1{d-2})$, decreases
up to $\max(\frac18,\frac1{d-2})$ and increases afterwards. Moreover
\begin{equation}\label{eq:exact}
 f_4(\tfrac18)=\tfrac14\log\tfrac{27}{8}=\tfrac L2 .
\end{equation}
\end{lemma}
\begin{proof}
$f_d'(r)=\frac{10r+1}{16r(1+r)}-\frac{d+8}{8(dr+1)}$, and
$(10r+1)(dr+1)-2(d+8)r(1+r)=(8d-16)r^2-(d+6)r+1=Q_d(r)$. For \eqref{eq:exact}:
$\frac9{16}\log\frac98+\frac1{16}\log\frac18-\frac38\log\frac38=\frac34\log3-\frac34\log2$.
\end{proof}

By Lemma~\ref{lem:onevar} the maximum of $f_d$ over an interval $(0,R]$ is attained at $R$ or at the first
critical point $\rho_d$ (if $\rho_d\le R$), where $\rho_2=\frac18$ and $\rho_d=\min(\frac18,\frac1{d-2})$ for
$d\ge3$. With the ranges of Lemma~\ref{lem:ranges} put
\[
 F_2:=\max_{(0,1]}f_2,\qquad F_d:=\max_{(0,3/d]}f_d,\qquad F^*_d:=\max_{(0,R^*_d]}f_d\qquad(d\ge3).
\]
All candidate points are rational, so these are explicit combinations of logarithms of rational numbers.

\begin{lemma}[numerical facts]\label{lem:num}
\textup{(a)} $F^*_4=\frac L2$, attained only at $r=\frac18$, and $F^*_d<\frac L2-0.0014$ for $3\le d\le119$, $d\ne4$.
\textup{(b)} $F_2<0.230041$ and $F_d<0.304973$ for $3\le d\le119$.
\end{lemma}
\noindent Some values: $F_2=0.230040$, $F^*_3=0.293861$, $F^*_4=0.304099$, $F^*_5=0.302650$, $F^*_6=0.300921$,
$F^*_7=0.295802$, $\max_dF_d=F_5=0.304972$; see Appendix~\ref{app:num} for how they are certified.

\begin{lemma}\label{lem:edge}
For every edge $e$ of every $G\in\cC$, $E_e\le L$. Consequently $\Phi(G)\le mL$.
\end{lemma}
\begin{proof}
Let $a=d_u\le b=d_v$.

\emph{Case A: $3\le a\le b\le119$.} Each end has a neighbour of degree at least $3$, namely the other end. By
Lemmas \ref{lem:sum}(b), \ref{lem:ranges}(c) and \ref{lem:num}(a), $E_e\le F^*_a+F^*_b\le L$.

\emph{Case B: $a=2$ and $b\le119$.} Then $b\ge3$, and $E_e\le F_2+F_b<0.230041+0.304973<L$.

\emph{Case C: $b\ge120$.} Dropping the non-positive terms $-g'(\rb_w)r_e$ in the definition of $E_e$,
\[
 E_e\le\log(1+r_e)+H_a(\rb_u)+H_b(\rb_v),\qquad H_d(s):=-\tfrac1d\log\bigl(\tfrac1d+s\bigr)+\tfrac9{16}\,\frac{s}{1+s}+\tfrac1{16},
\]
since $s\,g'(s)=\frac9{16}\frac s{1+s}+\frac1{16}$. For $d\ge120$ and $0<s\le\frac3d$ we have
$H_d(s)\le\frac{\log d}{d}+\frac{27}{16(d+3)}+\frac1{16}$, which decreases in $d$ and is less than $0.116116$ at $d=120$. For $2\le d\le119$, $H_d'$ has the sign of
$-16s^2+(9d-32)s-7$: if $d\le5$ this is negative; if $d\ge6$ its larger root exceeds $\frac3d$, so on the range of
Lemma~\ref{lem:ranges} $H_d$ decreases and then possibly increases. Hence $\sup H_d$ is the larger of
$H_d(0^+)=\frac{\log d}d+\frac1{16}$ and the value at the right end of the range, and both are at most
$\frac{\log3}{3}+\frac1{16}<0.428705$ (Appendix~\ref{app:num}). Finally $r_e\le x_ux_v\le\frac32\cdot\frac4{120}=\frac1{20}$
by \eqref{eq:I} and Lemma~\ref{lem:ranges}. Altogether $E_e<0.048791+0.428705+0.116116<L=0.608197\ldots$

The last assertion follows from Lemma~\ref{lem:sum}(a).
\end{proof}

\begin{remark}\label{rem:eq}
Equality $E_e=L$ forces $a=b=4$ and $\rb_u=\rb_v=r_e=\frac18$, the values of the $4$-regular tree; so
$\Phi(G)=mL$ if and only if $G$ is $4$-regular. The margin in Case~A is $0.00145$, attained at $d=5$.
\end{remark}

\section{Proofs of the main results}\label{sec:proof}

\begin{proof}[Proof of Theorem~\ref{thm:main}]
By Section~\ref{sec:red} it suffices to show $\tau(G)\le\tht^{\,m}$ for $G\in\cC$. Let $G_N$ be the graph obtained
from $N$ disjoint copies of $G$ by identifying one fixed vertex of all copies. Then $G_N\in\cC$,
$m(G_N)=Nm$, and $\tau(G_N)=\tau(G)^N$ because the number of spanning trees is multiplicative over blocks.
Lemmas \ref{lem:compare} and \ref{lem:edge}, applied to $G_N$, give $N\log\tau(G)\le NmL+C_1$ with $C_1$ independent
of $N$. Dividing by $N$ and letting $N\to\infty$ gives $\log\tau(G)\le mL$.
\end{proof}

\begin{proof}[Proof of Proposition~\ref{prop:sharp}]
By \cite{Lyons} (Theorem~3.3 and Example~3.16 in the arXiv version; see also \cite{McKay,ChungYau}),
$|V(G_j)|^{-1}\log\tau(G_j)\to3\log\frac32=\log\frac{27}8$
for connected $4$-regular graphs $G_j$ whose girth tends to infinity; such graphs exist, and $m(G_j)=2|V(G_j)|$.
\end{proof}

Corollaries \ref{cor:sup}, \ref{cor:planar} and \ref{cor:current} were derived in the introduction; in
Corollary~\ref{cor:sup} the supremum is not attained because the inequality of Theorem~\ref{thm:main} is strict.

\section{Literature}\label{sec:lit}

We searched for earlier statements of Theorem~\ref{thm:main} and found none. The search covered the papers cited
here and the papers citing \cite{Kenyon}, \cite{Stoimenow} and \cite{CKP} that are indexed in OpenAlex, as well as
arXiv and the OEIS. It was limited to open sources: we could not consult MathOverflow, Google Scholar or
publishers' sites. In particular we did not have access to \cite{Stoimenow}, \cite{McKay}, \cite{ChungYau},
\cite{Grimmett}, \cite{Kostochka} or the journal version of \cite{Kenyon}; what we say about them is taken from
\cite{CKP,GJKLZ} (Stoimenow), from \cite{GIKM,Lyons} (McKay, Chung--Yau, Kostochka), from \cite{ShrockWu}
(Grimmett) and from the arXiv version of \cite{Kenyon}. We therefore state plainly that the result may be known to
specialists, and we would be grateful for references.

What we found. For regular graphs the sharp constant is McKay's \cite{McKay} (see \cite{GIKM} for the statement
and \cite{ChungYau} for sharpness); for simple regular graphs the exponential rate $\tht$ already follows from it,
since the per-edge constants $1.7472$, $1.8371$, $1.8098$, $1.7571$ for $d=3,4,5,6$ are largest at $d=4$ (and they
decrease for $d\ge4$). For general
degree sequences we only found the bounds of Grimmett \cite{Grimmett} and Kostochka \cite{Kostochka}, whose
exponential rates per edge, $e^{2/e}=2.087$ and $3^{2/3}=2.080$, exceed $2$, and the unpublished bound of Young
mentioned by Kenyon. Lyons \cite{Lyons} shows that among transitive graphs
of a fixed degree the regular tree maximises tree entropy per vertex; we found nothing per edge or for irregular
covers. For planar graphs, see the introduction; recent work \cite{Burton,CK25} still treats the planar upper
bound as a conjecture. The formula of Lemma~\ref{lem:formula} is in \cite{BBGSS}, and the system \eqref{eq:I}
is the covariance system (2.5) of \cite{Watanabe} (with unit weights, at energy $0$).

\section{Questions}

\textup{(1)} Kenyon's first question: which graphs with $m$ edges maximise $\tau$? By exhaustive computation the
maximum of $\tau^{1/m}$ over connected multigraphs with at most $14$ edges is $1.667031$ ($\tau=1280$, $m=14$), far
from $\tht$: small graphs are not close to extremal.
\textup{(2)} The planar supremum. Since $\tht$ is exact for general graphs, a planar improvement must use faces,
that is, short cycles of the graph and of its dual. \textup{(3)} Is $\tau(G)\le\tht^{\,m}$ true with a
polynomial saving in $m$, as in McKay's bound for regular graphs?

\appendix
\section{Computer-assisted parts and verification package}\label{app:num}

The proof uses computation in one place only: the finitely many inequalities of Lemma~\ref{lem:num} and the three
constants of Case~C. The package accompanying this paper (\url{https://github.com/kosza123/flow-certificate-thresholds}, directory \code{spanning-trees-per-edge}, and the Zenodo record of this paper) contains the following.
\begin{itemize}
\item \code{verify/certify\_intervals.py}: evaluates $f_d$ at the rational candidate points for $2\le d\le119$ with
  $60$-digit interval arithmetic (\code{mpmath.iv}) and checks Lemma~\ref{lem:num}, the constants of Case~C, the
  inequality $k\le\tht^{\,k-1}$ in exact integers, and the constants $z_0,\alpha,\beta$ as exact rationals.
  
\item \begin{sloppypar}\code{lean/}: the Lean~4 package \code{TreeEntropy} (Lean 4.37.0-rc1, Mathlib commit
  \code{53dad9ed}; about $1600$ lines in eight modules), described in the paragraph on the status of
  verification. Its main statements are the theorems \code{tree\_entropy\_bound} (the edge inequality from the
  identities \eqref{eq:I} and the ranges) and \code{tree\_entropy\_bound\_of\_first\_step} (the same conclusion
  from the first-step relations) in the namespace \code{TreeEntropy}. The script \code{scripts/check.sh} builds the
  package, scans for placeholders, prints the axioms of $44$ theorems, runs positive and negative controls and
  replays every module in the kernel checker; a recorded clean-room run is included
  (\code{VERIFICATION\_CLEANROOM.txt}). The formal proof of Case~C does not use the sign analysis of $H_d'$ given
  in Section~\ref{sec:edge} and is in this sense independent of it.\end{sloppypar}
\item \code{verify/check\_class.py}: for random multigraphs of $\cC$ (including vertices of degree $2$) it solves
  \eqref{eq:h}, checks \eqref{eq:I}, the ranges of Lemma~\ref{lem:ranges}, the closed formula of
  Lemma~\ref{lem:formula} against a numerical integration of the return generating function, and
  $E_e\le L$. This is a consistency check, not part of the proof.
\item \code{verify/exhaustive.py}: checks $\tau(G)^2\,8^m<27^m$ in exact integer arithmetic for all connected
  loopless multigraphs with at most a given number of edges that are not trees (for a tree $\tau=1$), generated
  with \code{nauty} \cite{nauty}. Up to $14$ edges these are $8{,}350{,}006$ graphs, and there is no violation. This
  is a sanity check far from the sharp regime.
\end{itemize}

\section*{Author contributions and use of AI}
This work was carried out with substantial assistance from AI systems working under the author's direction. The
SAPIO MORPHEUS project uses models from OpenAI and Anthropic.
\begin{itemize}
\item The statement, the proof, the computations, the literature comparison and this manuscript were produced by AI
  agents based on Anthropic's Claude.
\item Separate AI agents of the same project, working with their own code, refereed the proof and then this text
  line by line and recomputed the numerical claims, including the exhaustive computation; another such agent
  carried out the literature search.
  
\item The Lean formalisation was written by a further AI agent based on Anthropic's Claude.
\item The author selected the direction, directed the work and made the publication decisions, and takes
  responsibility for the content.
\end{itemize}
In line with common publication policies, AI systems are not listed as authors. No external expert has reviewed
this manuscript at the time of writing.

\begin{thebibliography}{99}
\bibitem{BBGSS} J.~Banks, J.~Breuer, J.~Garza-Vargas, E.~Seelig, B.~Simon, \emph{A useful formula for periodic
  Jacobi matrices on trees}, Proc.\ Natl.\ Acad.\ Sci.\ USA \textbf{121} (2024), no.~23, e2315218121; arXiv:2309.00437.
\bibitem{Burton} S.~D.~Burton, \emph{The spectra of volume and determinant densities of links}, Topology Appl.\
  \textbf{211} (2016), 38--55; arXiv:1507.01954.
\bibitem{CK25} A.~Champanerkar, I.~Kofman, \emph{Geometric bounds for spanning tree entropy of planar lattice
  graphs}, arXiv:2505.05688 (2025).
\bibitem{CKP} A.~Champanerkar, I.~Kofman, J.~S.~Purcell, \emph{Geometrically and diagrammatically maximal knots},
  J.\ London Math.\ Soc.\ (2) \textbf{94} (2016), 883--908; arXiv:1411.7915v4.
\bibitem{ChungYau} F.~Chung, S.-T.~Yau, \emph{Coverings, heat kernels and spanning trees}, Electron.\ J.\ Combin.\
  \textbf{6} (1999), Research Paper~12.
\bibitem{Fomin} D.~Fomin, \emph{Elementary proof for the bounds of the complexity of a planar multigraph and the
  size of a prime rectangular squaring}, arXiv:2103.10523 (2021).
\bibitem{GJKLZ} J.~Ge, X.~Jin, L.~H.~Kauffman, P.~Lopes, L.~Zhang, \emph{Minimal sufficient sets of colors and
  minimum number of colors}, J.\ Knot Theory Ramifications \textbf{26} (2017), no.~9, 1743008; arXiv:1501.02421.
\bibitem{GIKM} C.~Greenhill, M.~Isaev, M.~Kwan, B.~D.~McKay, \emph{The average number of spanning trees in sparse
  graphs with given degrees}, European J.\ Combin.\ \textbf{63} (2017), 6--25; arXiv:1606.01586.
\bibitem{Grimmett} G.~R.~Grimmett, \emph{An upper bound for the number of spanning trees of a graph}, Discrete
  Math.\ \textbf{16} (1976), 323--324.
\bibitem{Kenyon} R.~Kenyon, \emph{Tiling a rectangle with the fewest squares}, J.\ Combin.\ Theory Ser.\ A
  \textbf{76} (1996), 272--291; arXiv:math/9411215.
\bibitem{Kostochka} A.~V.~Kostochka, \emph{The number of spanning trees in graphs with a given degree sequence},
  Random Structures Algorithms \textbf{6} (1995), 269--274.
\bibitem{Lyons} R.~Lyons, \emph{Asymptotic enumeration of spanning trees}, Combin.\ Probab.\ Comput.\ \textbf{14}
  (2005), 491--522; arXiv:math/0212165.
\bibitem{M2} F.~Majchrzak, \emph{Two-port substitution and the growth of the smallest current in planar networks}
  (2026), \url{https://github.com/kosza123/flow-certificate-thresholds}, directory \code{growth-constant}.
\bibitem{McKay} B.~D.~McKay, \emph{Spanning trees in regular graphs}, European J.\ Combin.\ \textbf{4} (1983),
  149--160.
\bibitem{nauty} B.~D.~McKay, A.~Piperno, \emph{Practical graph isomorphism, II}, J.\ Symbolic Comput.\ \textbf{60}
  (2014), 94--112.
\bibitem{ShrockWu} R.~Shrock, F.~Y.~Wu, \emph{Spanning trees on graphs and lattices in $d$ dimensions}, J.\ Phys.\ A
  \textbf{33} (2000), 3881--3902; arXiv:cond-mat/0004341.
\bibitem{Stoimenow} A.~Stoimenow, \emph{Maximal determinant knots}, Tokyo J.\ Math.\ \textbf{30} (2007), 73--97.
\bibitem{Watanabe} Y.~Watanabe, \emph{Gaussian belief propagation and periodic Jacobi matrices on trees: spectral
  curve, antibound poles, and fixed points}, arXiv:2611.29381 (2026).
\end{thebibliography}

\end{document}
